% Korean edition of sections/s2_4_feed.tex.
\seclabel
\claim{\secnum{2.4}{sec:feed}\quad 실시간 피드는 요청할 때만 답하고, 클라이언트가 \code{t}를 정확히 되돌려 보낼 때만 델타를 보낸다}

\lead{실시간 피드에 연결하려면 4단계를 거친다(\figref{signalr}).} 잔고, 주문, 섹터 비중은 \code{/hubs/main}의 ASP.NET Core SignalR 허브로 온다. HTTP 협상, WebSocket 업그레이드, 핸드셰이크를 마치면 클라이언트가 피드 요청을 보내고 서버가 하나씩 답한다.

% Korean edition of figs/signalr.tex.
% Time-space diagram in the style of Kurose & Ross Fig 2.7/3.39: two timelines, time runs down, sloped arrows,
% message names along the arrows, events beside the timelines. The time axis is broken twice (zigzags).
\fig[wide]{signalr}{협상과 핸드셰이크를 마친 뒤 서버는 피드 요청에만 답한다}[클라이언트와 실시간 피드 허브 \code{/hubs/main} 사이의 메시지. 시간은 아래로 흐른다. 연결 설정은 밀리초 단위, 그다음은 3\,s 폴링 주기]{%
\begin{tikzpicture}[x=1mm,y=1mm,
  f/.style={font=\fontsize{7}{8.4}\selectfont},
  host/.style={font=\fontsize{7.5}{9}\selectfont\bfseries,text=navy,anchor=south,align=center},
  tl/.style={greyM,line width=0.6pt},
  msg/.style={-{Stealth[length=1.6mm]},line width=0.7pt,navy},
  rep/.style={-{Stealth[length=1.6mm]},line width=0.7pt,greyD},
  ml/.style={font=\fontsize{7.5}{9}\selectfont,sloped,above,inner sep=0.7pt,text=navy},
  rl/.style={font=\fontsize{7.5}{9}\selectfont,sloped,below,inner sep=0.7pt,text=greyD},
  ev/.style={f,anchor=east,align=right,text=greyD},
  brk/.style={line width=0.6pt,greyM,fill=white},
  stp/.style={f,anchor=west,text width=48mm,align=left,text=greyD,execute at begin node={\hyphenpenalty=10000\exhyphenpenalty=10000\relax}},
  brace/.style={decorate,decoration={brace,amplitude=1mm},line width=0.5pt,navy}]
\def\C{0}\def\S{92}
% ---------------------------------------------------------------- hosts and timelines
\node[host] at (\C,1) {클라이언트\\(브라우저 또는 래퍼)};
\node[host] at (\S,1) {서버\\(SignalR 허브)};
\draw[tl] (\C,0) -- (\C,-106); \draw[tl] (\S,0) -- (\S,-106);
\draw[tl,-{Stealth[length=1.4mm]}] (\C,-106) -- (\C,-109); \draw[tl,-{Stealth[length=1.4mm]}] (\S,-106) -- (\S,-109);
\node[f,text=greyD,anchor=north] at (\C,-109.3) {시간};\node[f,text=greyD,anchor=north] at (\S,-109.3) {시간};
% broken-axis marks: two zigzags across both timelines
\newcommand{\sigzig}[1]{\foreach \x in {\C,\S}{\fill[white] (\x-1.6,#1+0.9) rectangle (\x+1.6,#1-0.9);
  \draw[greyM,line width=0.5pt] (\x-1.6,#1+1.4) -- (\x+1.6,#1+0.4); \draw[greyM,line width=0.5pt] (\x-1.6,#1-0.4) -- (\x+1.6,#1-1.4);}}
% ---------------------------------------------------------------- set-up (milliseconds)
\draw[msg] (\C,-6) -- node[ml]{\code{POST /hubs/main/negotiate?negotiateVersion=1}} (\S,-10);
\draw[rep] (\S,-11) -- node[rl]{200: \code{connectionId}, \code{connectionToken}, 전송 방식 3개} (\C,-15);
\draw[msg] (\C,-23) -- node[ml]{\code{GET /hubs/main?id=...}, \code{Upgrade: websocket}} (\S,-27);
\draw[rep] (\S,-28) -- node[rl]{\code{101 Switching Protocols}} (\C,-32);
\draw[msg] (\C,-40) -- node[ml]{\code{{"protocol":"json","version":1}} 0x1E} (\S,-44);
\draw[rep] (\S,-45) -- node[rl]{\code{{}} 0x1E} (\C,-49);
% events beside the client timeline
\node[ev] at (\C-2,-8) {연결 1, HTTP/2,\\REST용으로 이미 열림};
\node[ev] at (\C-2,-25) {연결 2: 새 TLS\\연결, HTTP/1.1};
\node[ev] at (\C-2,-42) {핸드셰이크};
% ---------------------------------------------------------------- poll cycle
\draw[msg] (\C,-56) -- node[ml]{\code{RequestPortfolio} \code{{intv:3, auto:true, lastT:0, d, pfId}}} (\S,-60);
\draw[rep] (\S,-61) -- node[rl]{\code{ReceivePortfolio}: 스냅샷 (\code{bs}, \code{ords}, \code{pos}, \code{msecs})} (\C,-65);
\draw[brace] (\C-2,-65) -- (\C-2,-56) node[midway,ev,xshift=-1.2mm] {응답\\20\textasciitilde 84\ms};
% 3 s wait with a break
\fill[navyM] (\C-0.6,-65.5) rectangle (\C+0.6,-76.5);
\sigzig{-71}
\draw[brace] (\C-2,-76.5) -- (\C-2,-65.5) node[midway,ev,xshift=-1.2mm] {3\,s 대기\\(축척 아님)};
\draw[msg] (\C,-77) -- node[ml]{\code{RequestPortfolio} \code{{..., lastT = t of previous reply}}} (\S,-81);
\draw[rep] (\S,-82) -- node[rl]{\code{ReceivePortfolio}: 델타 (\code{pos_diff}, \code{ords_diff}, \code{msecs_diff})} (\C,-86);
\fill[navyM] (\C-0.6,-86.5) rectangle (\C+0.6,-93);
\sigzig{-90}
\node[ev] at (\C-2,-90) {3\,s 대기 후 반복};
% ping: both directions, drawn level (schematic)
\draw[{Stealth[length=1.6mm]}-{Stealth[length=1.6mm]},line width=0.6pt,greyD,dashed] (\C,-100) -- node[f,above,inner sep=0.8pt,text=greyD]{유형 6 핑, 양쪽 모두 보냄. 서버는 15\,s마다} (\S,-100);
\node[ev] at (\C-2,-100) {유휴 연결을\\열어 둠};
% ---------------------------------------------------------------- step brackets on the right
\foreach \a/\b/\n/\t in {-5/-16/1/{\textbf{협상}: REST 연결에서 Bearer 토큰을 실은 HTTP/2 POST. 응답은 WebSocket, Server-Sent Events, 롱 폴링을 제시},
  -22/-33/2/{\textbf{연결}: 서버가 HTTP/2 위의 WebSocket[1]을 지원하지 않아 업그레이드가 두 번째 TLS 연결을 연다. 래퍼는 토큰을 헤더로 보내고, 브라우저는 URL에 넣는다},
  -39/-50/3/{\textbf{핸드셰이크}: 양쪽이 JSON 메시지에 합의하며, 메시지는 각각 0x1E 바이트로 끝난다},
  -55/-95/4/{\textbf{폴링}: 서버는 요청에 대한 답으로만 포트폴리오를 보낸다. 클라이언트는 응답마다 3\,s 뒤 다시 요청한다. \code{lastT}가 직전 응답의 \code{t}와 같으면 델타가 온다(\figref{lastt})}} {
  \draw[greyM,line width=0.5pt] (\S+3,\a) -- (\S+4,\a) -- (\S+4,\b) -- (\S+3,\b);
  \node[stp] at (\S+5.5,{(\a+\b)/2}) {\n.\ \t};
}
\end{tikzpicture}
}{남색 화살표는 클라이언트가, 회색 화살표는 서버가 보낸다. 시간은 아래로 흐른다. 연결 설정 부분은 측정한 지속 시간 없이 밀리초 단위로 그렸고, 지그재그 두 개는 줄여 그린 3\,s 대기를 뜻한다. 응답 시간 20\textasciitilde 84\ms{}는 새벽에 \code{auto}를 true와 false로 각각 두고 쟀다. 양쪽이 서로 독립적으로 핑을 보내므로 핑 화살표는 수평으로 그렸다. 유형 1은 상대편 메서드 호출, 유형 6은 핑, 유형 7은 종료다. 2026년 10월 28일 브라우저 기록에서 클라이언트는 직전 응답 3\,s 뒤에 다음 요청을 보냈다.}{브라우저 개발자 도구, Python 구현, DART Quant Team}


\lead{실시간 피드는 장 전이든 장중이든 요청하지 않으면 아무것도 보내지 않는다.} 브라우저 기록에서 브라우저는 응답을 받은 지 3\,s 뒤에 다시 요청한다. 새벽에 우리 클라이언트로 요청 플래그 \code{auto}를 true와 false로 각각 보냈을 때도(브라우저는 true를 보낸다) 같은 동작이었고, 응답은 20\textasciitilde 84\ms{} 안에 왔다. 가격이 움직이는 09:05\textasciitilde 09:08에 피드 요청 하나를 보내고 60\,s 동안 가만히 있던 클라이언트는 15\,s마다 오는 SignalR keep-alive 핑(메시지 유형 6) 외에는 아무것도 받지 않았다.

\lead{같은 연결에서 직전 응답의 \code{t}를 되돌려 보내면 바뀐 것만 온다.} \code{lastT=0}으로 요청하면 잔고(\code{bs}), 주문(\code{ords}), 포지션(\code{pos}), 시장 섹터 비중(\code{msecs})을 담은 스냅샷이 오며, \ref{sec:guidelines}절에서 대회 가이드라인에 쓰는 작은 최상위 필드 5개도 함께 온다. 직전 응답의 \code{t}를 \code{lastT}로 실은 요청은 델타를 받는다. 델타에서는 위 4개 필드가 \code{null}이고 대신 \code{pos_diff}, \code{ords_diff}, \code{msecs_diff}가 나타나며, 각각 \code{added}, \code{removed}, \code{updated}, \code{ords_chgd} 필드를 가진다. 여기에 바뀐 잔고 필드를 담은 딕셔너리 \code{bs_updated}가 붙는다. 우리 주문(\ref{sec:orders}절) 직후 응답 하나는 인자가 \code{null}이었다. 이를 널 응답이라 부르며, 델타 필드는 있지만 변경이 없는 빈 델타와 구분한다. 수신 루프는 널 응답을 받아들이고 다음에는 스냅샷을 요청해야 한다.

% Korean edition of figs/replysizes.tex.
\fig{replysizes}{델타는 개장 전 스냅샷의 1/17, 장중에는 약 1/8이다}[응답 유형과 시간대별 \code{ReceivePortfolio} 응답 하나의 크기(바이트). 2026년 11월 24일, 3\,s마다 폴링]{%
{\tabfont\begin{tabularx}{\linewidth}{>{\RaggedRight\arraybackslash}X R{26mm} R{50mm} R{14mm}}
\thd{시점} & \thd{스냅샷 (\code{lastT=0})} & \thd{델타 (\code{t} 정확히 되돌림)} & \thd{비율}\\
\rowsep
개장 전 (09:00 이전) & 9,915 & 577 (빈 델타) & 17\\
\rowsep
장중 & 11,344 & 1,053\textasciitilde 1,912 (평균 약 1,400) & 약 8\\
\end{tabularx}}
}{스냅샷은 잔고(\code{bs}), 주문(\code{ords}), 포지션(\code{pos}), 시장 섹터 비중(\code{msecs}) 전체와 작은 최상위 필드 5개를 담는다. 델타에서는 그 4개 필드가 \code{null}이고, \code{pos_diff}, \code{ords_diff}, \code{msecs_diff}가 추가·삭제·변경된 레코드를, \code{bs_updated}가 바뀐 잔고 필드를 싣는다. 개장 전에는 아무것도 바뀌지 않으므로 델타가 비어 있다. 장중에는 \code{pos_diff.updated}가 응답마다 보유 종목 23개 중 8\textasciitilde 23개를 담았고, 이것이 델타 크기의 범위를 정한다. 비율 = 스냅샷 크기 ÷ 델타 크기(장중은 평균 델타).}{실시간 피드에 연결한 Python 클라이언트, DART Quant Team}


\lead{장중 델타는 스냅샷의 약 8분의 1 크기다(\figref{replysizes}).} 개장 전 빈 델타는 스냅샷의 17분의 1이었다. 가격이 움직이는 동안 델타는 평균 약 1,400바이트, 스냅샷은 11,344바이트였다. \code{pos_diff.updated}는 응답마다 보유 종목 23개 중 8\textasciitilde 23개를 담았고, \code{msecs_diff}는 대략 4번째 응답마다(약 12\,s) 바뀌었다. \code{t}는 응답마다 3씩, 가끔 4씩 늘었으므로 벽시계 초를 따른다.

{\RaggedRight \lead{\code{updated}는 각 레코드에서 바뀐 필드만 담으며, 키는 레코드 \code{Id}를 문자열로 쓴 것이다.} 한 응답의 두 항목은 \code{"5745534": {"close": 89900, "prft": 1008000}}와 \code{"5745535": {"close": 1180000, "prft": 3230000, "chPct": 8.76, "wei": 4.86}}였다. 스냅샷의 포지션 레코드가 가진 키 28개 중 2\textasciitilde 4개다(부록~\ref{app:glossary}에 28개 모두 있다). \code{ords_diff}도 형태가 같다. 예: \code{{"added": [], "removed": [], "updated": {"1649882": {"cumQty": 6, "avgPx": 352416.67}}, "ords_chgd": {}}}. 스냅샷 레코드는 \code{Id}를 정수로 담으므로, 델타 키를 그대로 쓰면 어떤 레코드와도 맞지 않는다. 래퍼는 병합 전에 각 키를 \code{int()}로 바꾼다. 키는 포지션 레코드 ID이고 종목코드(\code{prodId})는 레코드의 별도 필드이므로, 래퍼는 \code{prodId}에서 \code{Id}로 가는 색인도 따로 둔다.\par}

\lead{이렇게 델타를 병합하자 15회 중 15회 모두 서버의 포지션과 주문이 필드 단위로 그대로 재현됐다.} 09:35에 연결 2개가 15회에 걸쳐 같은 순간에 실시간 피드를 폴링했다. 연결 A는 \code{updated} 필드를 \code{int(Id)}로 병합하며 델타를 적용했고, 연결 B는 매번 전체 스냅샷을 요청했다(\code{lastT=0}). 15회 모두 두 연결이 같은 \code{t}를 받았고, A에 누적된 \code{pos}와 \code{ords} 레코드는 B의 스냅샷과 필드 단위로 같았다. 15회 중 15회 모두 차이가 0건이었다. 이 검사는 \code{pos}와 \code{ords}만 비교했고, 잔고와 \code{msecs}는 필드 단위로 비교하지 않았다. 잔고는 \code{bs_updated}로 최신 상태를 유지한다. 10:33에 델타 20개 중 20개에서 \code{bs_updated}가 딕셔너리로 왔고, 키는 \code{prft}, \code{na}, \code{rt}, \code{nav}, \code{stkAmt}였으며 가끔 \code{pfStd}, \code{rank}도 있었다. 래퍼는 이를 \code{bs}에 병합한다(\ref{sec:oneshare}절). 앞선 검사에서는 병합 상태를 마지막 델타보다 약 3\,s 뒤에 받은 스냅샷과 비교했는데, 레코드 13개가 \code{close}, \code{prft}, \code{wei}, \code{chPct}에서만 달랐다. 3\,s 동안 가격을 따라 움직이는 필드다. \code{t}가 같은 응답끼리 비교하면 이 차이가 사라진다.

\lead{후보 규칙 3개 중 응답 11개에 모두 맞는 것은 (a)뿐이다. 서버는 \code{lastT}가 이 연결에서 마지막으로 보낸 \code{t}와 같을 때만 델타를 보낸다.} 후보 규칙은 다음과 같다. (a) \code{lastT}가 이 연결에서 마지막으로 보낸 \code{t}와 같다. (b) \code{lastT}가 최근 시각이기만 하면 된다. (c) \code{lastT}가 어느 연결이든 이 계정에 마지막으로 보낸 \code{t}와 같다. 개장 전인 06:06에 연결 2개에서 요청 11개의 \code{lastT}를 바꿔 가며, 규칙마다 다른 응답을 예측하도록 했다(\figref{lastt}). 이 테스트의 델타는 모두 비어 있었다. 이후 09:35 병합 검사는 델타에 변경이 실리는 장중에 진행됐고, 매번 \code{t}를 되돌려 보낸 연결 A는 15회 모두 스냅샷과 일치하는 델타를 계속 받았다. 이 결과는 (a)와 맞지만 (b)와 (c)를 시험하지는 않는다. 정확히 되돌려 보내면 세 규칙 모두 델타를 예측하기 때문이다.

% Korean edition of figs/lastt.tex.
\fig[wide]{lastt}{모든 응답을 예측하는 것은 규칙 (a)뿐이다. (b)는 5단계에서, (c)는 1단계에서 틀린다}[보낸 \code{lastT}별 예측 응답 유형과 관찰 결과. 연결 2개에서 3\,s 간격 요청 11개, 2026년 11월 24일 06:06 KST]{%
{\tabfont\begin{tabularx}{\linewidth}{R{6mm} >{\RaggedRight\arraybackslash}X >{\centering\arraybackslash}p{9mm} >{\centering\arraybackslash}p{17mm} >{\centering\arraybackslash}p{17mm} >{\centering\arraybackslash}p{17mm} >{\centering\arraybackslash}p{17mm}}
\thd{단계} & \thd{보낸 \code{lastT}} & \thd{연결} & \thd{(a) 이 연결의 마지막 \code{t}} & \thd{(b) 아무 최근 시각} & \thd{(c) 이 계정의 마지막 \code{t}} & \thd{관찰}\\
\rowsep
1 & 0 & A & \textcolor{greyD}{전체} & \textcolor{greyD}{전체} & \textcolor{greyD}{전체} & \textbf{전체}\\
\rowsep
2 & 직전 응답의 \code{t} & A & \textcolor{greyD}{델타} & \textcolor{greyD}{델타} & \textcolor{greyD}{델타} & \cellcolor{navyXL}\textbf{델타}\\
\rowsep
3 & 그 전 응답의 \code{t} (6초 전) & A & \textcolor{greyD}{전체} & \cellcolor{navyL}\textbf{델타\,$\times$} & \textcolor{greyD}{전체} & \textbf{전체}\\
\rowsep
4 & 직전 \code{t} 빼기 1 (보낸 적 없는 값) & A & \textcolor{greyD}{전체} & \cellcolor{navyL}\textbf{델타\,$\times$} & \textcolor{greyD}{전체} & \textbf{전체}\\
\rowsep
5 & 직전 \code{t} 빼기 60 & A & \textcolor{greyD}{전체} & \cellcolor{navyL}\textbf{델타\,$\times$} & \textcolor{greyD}{전체} & \textbf{전체}\\
\rowsep
6 & 직전 \code{t} 빼기 3,600 & A & \textcolor{greyD}{전체} & \textcolor{greyD}{해당 없음} & \textcolor{greyD}{전체} & \textbf{전체}\\
\rowsep
7 & 직전 \code{t} 더하기 100 & A & \textcolor{greyD}{전체} & \textcolor{greyD}{해당 없음} & \textcolor{greyD}{전체} & \textbf{전체}\\
\rowsep
8 & 직전 응답의 \code{t} & A & \textcolor{greyD}{델타} & \textcolor{greyD}{델타} & \textcolor{greyD}{델타} & \cellcolor{navyXL}\textbf{델타}\\
\rowsep
8$'$ & 같은 값을 다시 & A & \textcolor{greyD}{전체} & \cellcolor{navyL}\textbf{델타\,$\times$} & \textcolor{greyD}{전체} & \textbf{전체}\\
\specialrule{0.8pt}{1pt}{1pt}
9 & 연결 A의 마지막 \code{t} & B & \textcolor{greyD}{전체} & \cellcolor{navyL}\textbf{델타\,$\times$} & \cellcolor{navyL}\textbf{델타\,$\times$} & \textbf{전체}\\
\rowsep
10 & 직전 응답의 \code{t} & B & \textcolor{greyD}{델타} & \textcolor{greyD}{델타} & \textcolor{greyD}{델타} & \cellcolor{navyXL}\textbf{델타}\\
\end{tabularx}}
}{$\times$가 있는 중간 파랑 칸은 관찰과 어긋나는 예측이고, 회색 글씨는 관찰과 맞는 예측이다. 관찰 열에서 델타는 연한 회청색으로 칠했다. 해당 없음: 규칙 (b)는 1시간 전의 \code{t}(6단계)나 서버보다 100\,s 앞선 \code{t}(7단계)가 최근인지 말하지 않으므로 거기서는 어느 응답 유형도 예측하지 않는다. 3단계는 서버가 6\,s 전에 실제로 보낸 \code{t}를, 4단계는 서버가 보낸 적 없는 값을 보낸다. 5\textasciitilde 7단계는 허용 구간이 있는지 시험한다. 8$'$단계는 방금 쓴 값을 반복한다. 9단계는 연결 A의 마지막 \code{t}로 연결 B를 연다. 응답에 \code{pos}가 있으면 전체, \code{pos_diff}가 있으면 델타다. 개장 전 전체 응답은 9,915바이트, 델타는 577바이트였다(\figref{replysizes}).}{DART Quant Team}


\lead{서버는 연결마다 마지막으로 보낸 스냅샷과 그 \code{t}를 기억하는 것으로 보인다.} \code{lastT}가 그 \code{t}와 같으면 그 스냅샷 이후 바뀐 필드를 보내고, 아니면 전부 보낸다. 어느 쪽이든 현재 상태가 새 기준이 된다.

\lead{규칙 (a)에서는 델타를 건너뛰면 반드시 스냅샷으로 이어지므로 래퍼의 상태가 서버와 어긋나지 않는다.} 요청을 하나만 보내 두면 응답은 타임아웃으로만 사라질 수 있다. 래퍼는 요청마다 마감 시간을 20\,s로 한 번 정하고(서버의 15\,s 핑으로 연장되지 않는다), 넘기면 연결을 끊고 다시 연결한다. 새 연결은 \code{lastT=0}으로 시작하므로 다음 응답은 스냅샷이고, 클라이언트가 델타를 하나 건너뛴 채 다음 델타를 적용하는 일은 생기지 않는다. 첫 응답이 오기 전에 두 번째 요청을 보내면 그 \code{lastT}는 이미 낡았으므로 응답은 전체 스냅샷이 된다. 개장 전 기준으로 빈 델타의 17배 크기다(\figref{replysizes}).

\lead{\code{t}는 Unix 시간보다 정확히 9시간(32,400\,s) 앞서 있다.} 05:48:27 KST에 받은 응답의 \code{t}는 1795499307이었다. 이를 UTC로 읽으면 11월 24일 05:48:27로 응답의 \code{kstNow}와 맞지만, 05:48:27 KST의 Unix 시간은 1795466907이다. 서버가 한국 시각을 UTC인 것처럼 에포크 초로 바꾸는 것으로 보인다. 다른 시각은 올바르다. JWT의 \code{exp}는 실제 Unix 시간이고, 주문 레코드는 \code{Z}로 끝나는 UTC 문자열이다. 규칙 (a)에서는 클라이언트가 \code{t}로 계산하지 않고 되돌려 보내기만 하면 이 차이가 문제 되지 않는다.

\lead{열린 WebSocket은 토큰이 만료돼도 계속 동작하지만, 새 연결에는 유효한 토큰이 필요하다.} 07:21:41에 만료되는 토큰으로 연결을 열고 3\,s마다 폴링을 계속했다(\figref{token}). 응답은 만료 후 8분 동안 모두 382개가 계속 왔다. REST 요청이 거부되는 300\,s를 넘긴 시점이다. 그 뒤 같은 토큰으로 보낸 협상은 401을 받았다. 서버는 8분 동안 토큰을 다시 검사하지 않았는데, 이는 ASP.NET Core SignalR이 기본으로 연결할 때만 토큰을 검사한다는 것과 맞는다[3].

\lead{피드 오류가 나면 래퍼는 다시 연결하고, 피드가 26\,s 동안 조용하면 주문을 막는다.} 502, 401, 잘못된 JSON, 타임아웃은 연결을 닫는다. 래퍼는 0부터 min(30, 2\textsuperscript{\textit{n}})\,s 사이에서 균등하게 뽑은 무작위 시간만큼 기다린 뒤 다시 연결한다. \textit{n}은 연속 실패 횟수다(full-jitter 백오프). 래퍼는 마지막 응답이 도착한 시각을 기록하며, 그 응답이 2\,×\,3 + 20 = 26\,s(폴링 간격 2회와 응답 타임아웃)보다 오래되면 읽기, 주문, \code{wait()}이 \code{StaleFeed}를 던진다. 협상에서 받은 403이나 429는 예외로, 다시 연결하지 않고 피드를 멈춘다. 운영자가 이 신호를 어떻게 보고 처리하는지는 \ref{sec:design}장에 있다.

\lead{수신 루프는 이 결과에서 나온다(\figref{loop}).} 루프는 요청을 하나만 보내 두고, 스냅샷이면 상태를 교체하고, 델타면 바뀐 필드를 해당 레코드에, \code{bs_updated}를 \code{bs}에 병합한다. 널 응답 뒤에는 \code{lastT}를 0으로 되돌리고, \code{t}는 해석하지 않고 되돌려 보낸다. 다시 연결할 때는 토큰을 요청하고(때가 됐으면 이때 갱신된다, \ref{sec:auth}절), 다시 협상하고, 스냅샷부터 시작한다.

% Korean edition of figs/loop.tex.
% Pseudocode of Hub.stream / poll_once / assert_fresh in pytimefolio/hub.py, comments beside the code.
\fig[wide]{loop}{수신 루프는 요청을 하나만 보내 두고, 오류가 나면 다시 연결해 스냅샷부터 시작한다}[\code{hub.py}의 실시간 피드 루프 의사코드와 단계별 이유]{%
{\fontsize{7.6}{9.8}\selectfont
\newcommand{\cm}[1]{\leavevmode{\rmfamily\color{greyD}#1}}%
\newcommand{\I}[1]{\hspace*{#1mm}}%
\begin{tabular}{@{}>{\ttfamily\raggedright\arraybackslash}p{104mm} >{\RaggedRight\arraybackslash}p{70mm}@{}}
poll(hub):                                                   & \\
\I{4}n $\leftarrow$ 0; lastReply $\leftarrow$ never         & \cm{n은 연속 실패 횟수}\\
\I{4}loop:                                                   & \\
\I{8}try:                                                    & \\
\I{12}if not connected:                                      & \\
\I{16}negotiate; open WebSocket; handshake                   & \cm{토큰이 720\,s보다 오래됐거나 거부되면 먼저 갱신(\ref{sec:auth}절)}\\
\I{16}lastT $\leftarrow$ 0                                   & \cm{새 연결은 스냅샷부터 시작}\\
\I{12}send(RequestPortfolio(lastT, d = currBizDay, pfId))   & \\
\I{12}p $\leftarrow$ await receive(deadline = send + 20\,s)   & \cm{요청은 하나만. 핑은 마감을 늘리지 않음}\\
\I{12}if p = null: lastT $\leftarrow$ 0                      & \cm{널 응답: 상태 유지, 다음 응답은 스냅샷}\\
\I{12}elif p.pos $\neq$ null: state $\leftarrow$ p; lastT $\leftarrow$ p.t & \cm{스냅샷: 상태 교체}\\
\I{12}else: merge(state, p); lastT $\leftarrow$ p.t          & \cm{델타: 병합. \code{t}는 되돌리기만 하고 계산에 쓰지 않음}\\
\I{12}lastReply $\leftarrow$ now; n $\leftarrow$ 0           & \cm{아래 신선도 검사용으로 기록}\\
\I{8}except BlockedError:                                     & \cm{협상에서 403이나 429: 쓰기는 이미 중단됨}\\
\I{12}close WebSocket; raise                                 & \cm{피드 태스크 종료, 재연결 없음}\\
\I{8}except any other error:                                 & \cm{502, 401, 잘못된 JSON, 타임아웃, 닫힌 소켓}\\
\I{12}close WebSocket; n $\leftarrow$ n + 1                  & \\
\I{12}await sleep(uniform(0, min(30, 2$^{\mathrm{n}}$)))       & \cm{full-jitter 백오프, 최대 30\,s. 그다음 재연결}\\
\I{12}continue                                               & \\
\I{8}await sleep(3)                                          & \cm{응답마다 3\,s 뒤 다시 요청}\\[3pt]
merge(state, p):                                             & \\
\I{4}for each diff in p.pos\_diff, p.ords\_diff, p.msecs\_diff: & \\
\I{8}add diff.added; drop diff.removed                       & \\
\I{8}for (i, f) in diff.updated: state[int(i)] $\mid$= f     & \cm{바뀐 필드만. 키는 문자열로 된 레코드 \code{Id}}\\
\I{4}state.bs $\mid$= p.bs\_updated                          & \cm{순자산(\code{na})을 최신으로 유지}\\[3pt]
assertFresh():                                               & \cm{읽기, 주문, \code{wait()} 전에 호출}\\
\I{4}if now $-$ lastReply > 2 $\times$ 3 + 20 = 26\,s: raise StaleFeed & \cm{폴링 2회 누락과 20\,s 응답 타임아웃}\\
\end{tabular}}}{\code{updated}의 키는 문자열로 온 레코드 \code{Id}이고 스냅샷의 \code{Id}는 정수이므로 \code{int(i)}로 바꾼다. $\mid$=는 바뀐 필드의 딕셔너리를 저장된 레코드에 병합한다. 20\,s 마감은 요청마다 한 번 정하므로, 서버가 15\,s마다 보내는 핑이 잃어버린 응답을 계속 기다리게 만들 수 없다. \code{uniform(0, x)}는 0과 \code{x}초 사이에서 균등하게 대기 시간을 뽑고, 성공하면 n을 0으로 되돌린다. 토큰 갱신은 지연 방식이다. 피드는 연결하거나 다시 연결할 때만 토큰을 요청하며, 열린 WebSocket은 토큰이 만료돼도 계속 동작한다(\figref{token}). 두 번째 연결에서 받은 \code{t}가 같은 스냅샷과 15회 비교한 테스트에서 병합한 \code{pos}와 \code{ords} 레코드는 매회 필드 단위로 일치했다(\ref{sec:feed}절).}{\code{pytimefolio/hub.py} v0.1.0, DART Quant Team}

\verdict{수신 루프는 요청을 하나만 보내 두고 \code{t}를 해석하지 않고 되돌려 보낸다. 시각이 필요하면 래퍼는 응답의 \code{Date} 헤더에서 얻은 서버 시계를 쓴다(\ref{sec:auth}절). \code{kstNow}는 저장하지만 쓰지 않는다(D8). 델타는 \code{int(Id)}로, 바뀐 필드만 병합한다(D14). 토큰을 갱신해도 열린 WebSocket은 그대로 둔다. 다시 연결할 때는 먼저 토큰을 요청하고(때가 됐으면 갱신), 그 토큰을 헤더로 보낸다. 그래서 액세스 토큰이 URL에 들어가지 않으며, URL에는 수명이 짧은 연결 토큰만 실린다(D9).}
